\section{Measuring Similarity to Never-Training}
\label{sec:evaluation}

For each concept we compare a shared control $M_C$, its concept-excluded twin
$M_N$, and an erased model $M_E$ produced from $M_C$. The excluded effect is
$M_N-M_C$, the erasure effect is $M_E-M_C$, and the residual
$M_E-M_N$ measures mismatch. We apply this design to four complementary notions
of proximity rather than defining success with one aggregate score.

\subsection{Behavioral Similarity}

EMBER supplies 100 target questions and 100 neighboring-domain questions per
concept, each divided into equal selection and held-out halves. Ancient Rome's
neighboring set covers other historical civilizations; Baseball's covers other
sports. Because these base models do not reliably follow a generated-letter
multiple-choice format, questions are rewritten as declarative stems and each
answer option is scored as a continuation.

For option $a$ and stem $x$, we compute
\begin{equation}
s(a,x)=\frac{\log P(a\mid x)-\log P(a\mid \varnothing)}{|a|_{\mathrm{char}}}.
\end{equation}
The subtraction reduces preference for intrinsically common option strings and
the denominator mitigates option-length bias. Our primary per-question value is
the gold option's $s(a,x)$; models are compared on paired questions. Argmax
accuracy is secondary because 50 binary outcomes per split are too coarse to
resolve many changes visible in likelihood.

\subsection{Distributional Similarity}

Each model scores the same 3,000 blacklisted chunks and 5,004 non-blacklisted
chunks sampled to match the variable-sequence-length bucket distribution. For
model pair $(A,B)$, let $\Delta_H(A,B)$ be the mean paired increase in
per-token negative log-likelihood (NLL) on held-out concept chunks and
$\Delta_R(A,B)$ the corresponding increase on control chunks. We report
\begin{equation}
\mathrm{DiD}(A,B)=\Delta_H(A,B)-\Delta_R(A,B),
\end{equation}
which subtracts global model drift. We also inspect quantiles of the per-chunk
shift: equal means can conceal a uniform mild degradation versus a small tail
of severely damaged passages.

\placeholder{Add output-distribution KL divergence only if the team implements
it on a clearly specified prompt/token set. NLL difference-in-differences is
the completed distributional measurement.}

\subsection{Parameter-Space Similarity}

Flattening aligned floating-point parameters, define the erasure update
$d_E=W_E-W_C$ and target displacement $d_N=W_N-W_C$. We report
\begin{align}
\mathrm{cosine} &= \frac{d_E^\top d_N}{\lVert d_E\rVert\lVert d_N\rVert},\\
\mathrm{progress} &= \frac{d_E^\top d_N}{\lVert d_N\rVert^2},\\
\mathrm{residual} &= \frac{\lVert d_E-d_N\rVert}{\lVert d_N\rVert}.
\end{align}
Cosine measures direction, progress measures how much of the target direction
the edit traverses, and residual equals zero only at the excluded model. Global
measurements include all aligned parameters; embedding-only results isolate the
space EMBER can directly change. This analysis assumes matched training, making
the hardware qualification in Section~\ref{sec:experimental-design} relevant.

\subsection{Specificity and Capability Preservation}

Specificity is measured at three distances from the target: neighboring-domain
questions, length-matched general corpus chunks, and general OLMES tasks. A
desirable edit approaches $M_N$ on target measurements without moving farther
from it on controls. Neighboring questions are semantically demanding but
small; control-chunk NLL has much higher power but is not topic matched; general
benchmarks detect broad capability loss but need not be sensitive to either
target concept.

\subsection{Statistical Analysis}

Question and same-set chunk contrasts are paired by item. We report paired
standardized effects ($d_z$) and two-sided sign-flip permutation tests. For the
held-out-versus-control difference-in-differences, whose two sets are unpaired,
we permute set labels. \placeholder{Add 95\% paired bootstrap confidence
intervals to every headline contrast and state the resampling unit.} These
statistics quantify item-level uncertainty only; with one training seed per
condition they do not measure training-run variability. Erasure settings are
selected without consulting the final concept-excluded comparison wherever the
record permits; any exception must be labeled exploratory.
